\chapter{Statistics, Geometry, and Similarity}
\label{ch:foundations}
\label{ch:statistics}
\label{ch:linear-algebra}

\chapterguide%
  {\chapref{ch:data}, \chapref{ch:learning}, and comfort with arithmetic, fractions, powers, and basic algebra.}
  {use statistics, probability, vectors, matrices, distances, and gradients; recognize uncertainty and scaling effects.}

Statistics describes variation, uncertainty, and association; geometry describes vectors, distances, and transformations.

\begin{secondpass}
\textbf{Second pass:} standard errors, confidence intervals, bootstrap, distributions, Bayes' algebra, inverses, eigenvectors/SVD, and curvature. These need not delay the first models.
\end{secondpass}

\section*{Part A: Describe data with statistics and probability}
\section{Describe a distribution with center and spread}

A \term{distribution} describes which values occur and how often they occur. Two questions come first:

\begin{enumerate}
  \item Where is the center?
  \item How spread out are the values?
\end{enumerate}

\subsection{Mean and median}

For values $x_1,\ldots,x_N$, the population mean is
\[
  \mu=\frac{1}{N}\sum_{i=1}^{N}x_i.
\]
The symbol $\sum$ means summation: add the $N$ values, then divide by $N$.

The sample mean is usually written $\bar{x}$ and uses the same calculation on a sample.

The \term{median} is the middle value after sorting. It is less affected by extreme values.

\begin{workedexample}
For $30,32,35,38,500$, the mean is $127$ and median $35$. The median better describes typical income here because the extreme value pulls the mean upward.
\end{workedexample}

\subsection{Variance and standard deviation}

Variance measures average squared distance from the mean:
\[
  \sigma^2=\frac{1}{N}\sum_{i=1}^{N}(x_i-\mu)^2.
\]
The standard deviation is
\[
  \sigma=\sqrt{\sigma^2}.
\]
Standard deviation is easier to interpret because it uses the original unit.

\begin{workedexample}
For $2,4,6$, the mean is $4$. The distances from the mean are $-2,0,2$. Squaring gives $4,0,4$, so the population variance is $8/3\approx2.67$ and the standard deviation is about $1.63$.
\end{workedexample}

\subsection{Quantiles and the interquartile range}

A quantile is a cut point in sorted data. The first quartile $Q_1$ has about 25\% of values below it, and the third quartile $Q_3$ has about 75\% below it. The interquartile range is
\[
  \operatorname{IQR}=Q_3-Q_1.
\]
It measures the width of the middle half of the data and is resistant to extreme values.

\begin{keyidea}
Use the mean and standard deviation for roughly symmetric data without severe outliers. Use the median and IQR when the distribution is skewed or contains extreme values.
\end{keyidea}

\section{Mean versus median, drawn}

The mean and median agree on symmetric data and disagree sharply on skewed data. That disagreement is itself a useful diagnostic.

\begin{mlfigure}
\begin{tikzpicture}[baseline=(current bounding box.north)]
\begin{axis}[
  width=0.46\linewidth, height=4.2cm, axis lines=left,
  title style={font=\small\bfseries\color{MLNavy}}, title={Symmetric: they agree},
  xlabel={Value}, ylabel={How many},
  label style={font=\scriptsize\color{MLMuted}}, tick label style={font=\scriptsize\color{MLMuted}},
  ybar, bar width=7pt, ymin=0, ymax=20, xmin=0, xmax=9
]
\addplot[fill=MLPaleBlue,draw=MLBlue] coordinates {(1,2)(2,6)(3,12)(4,17)(5,12)(6,6)(7,2)};
\draw[MLRed,very thick,dashed] (axis cs:4,0) -- (axis cs:4,19);
\node[font=\scriptsize\color{MLRed}] at (axis cs:6.5,17) {mean = median};
\end{axis}
\end{tikzpicture}
\hfill
\begin{tikzpicture}[baseline=(current bounding box.north)]
\begin{axis}[
  width=0.46\linewidth, height=4.2cm, axis lines=left,
  title style={font=\small\bfseries\color{MLNavy}}, title={Skewed: they split apart},
  xlabel={Value}, ylabel={How many},
  label style={font=\scriptsize\color{MLMuted}}, tick label style={font=\scriptsize\color{MLMuted}},
  ybar, bar width=7pt, ymin=0, ymax=20, xmin=0, xmax=9
]
\addplot[fill=MLPaleOrange,draw=MLOrange] coordinates {(1,18)(2,13)(3,8)(4,5)(5,3)(6,2)(7,1)};
\draw[MLTeal,very thick,dashed] (axis cs:2,0) -- (axis cs:2,19);
\draw[MLRed,very thick,dashed] (axis cs:2.8,0) -- (axis cs:2.8,19);
\node[font=\scriptsize\color{MLTeal}] at (axis cs:5.6,17) {median};
\node[font=\scriptsize\color{MLRed}] at (axis cs:5.6,13.5) {mean (pulled right)};
\end{axis}
\end{tikzpicture}
\end{mlfigure}

\begin{keyidea}
A mean far above the median suggests right skew. Consider a log transform or robust scaler; extreme errors also affect RMSE more than MAE (\chapref{ch:evaluation}).
\end{keyidea}

\section{The normal distribution}

The normal distribution is the familiar bell-shaped curve. It is controlled by a mean $\mu$ and standard deviation $\sigma$:
\[
  f(x)=\frac{1}{\sigma\sqrt{2\pi}}
  \exp\!\left[-\frac{1}{2}\left(\frac{x-\mu}{\sigma}\right)^2\right].
\]

\begin{mlfigure}
\begin{tikzpicture}
\begin{axis}[
  width=0.82\linewidth,height=5.3cm,axis lines=left,
  xlabel={value},ylabel={density},ytick=\empty,
  xtick={-2,-1,0,1,2},
  xticklabels={$\mu-2\sigma$,$\mu-\sigma$,$\mu$,$\mu+\sigma$,$\mu+2\sigma$},
  domain=-3.2:3.2,samples=140,ymin=0,ymax=0.45,clip=false
]
\addplot[very thick,MLBlue] {exp(-x^2/2)/sqrt(2*pi)};
\draw[dashed,MLMuted] (axis cs:0,0)--(axis cs:0,0.399);
\draw[{Stealth}-{Stealth},thick,MLTeal] (axis cs:-1,0.24)--(axis cs:1,0.24)
 node[midway,above,font=\small] {about 68\%};
\end{axis}
\end{tikzpicture}
\end{mlfigure}

Many measurement errors and averages are approximately normal, but real data should not be assumed normal without inspection.

\begin{intuition}
$\mu$ sets the bell's center; $\sigma$ sets its width.
\end{intuition}

\section{Population, sample, and sampling bias}

A \term{sample} is the observed subset of a \term{population}. Sample statistics estimate population parameters.

For a sample of size $n$, the common sample variance is
\[
  s^2=\frac{1}{n-1}\sum_{i=1}^{n}(x_i-\bar{x})^2.
\]
The denominator $n-1$ corrects the tendency to underestimate population variance when deviations are measured around the sample mean.

\begin{pitfall}
More observations do not correct systematic sampling bias.
\end{pitfall}

\section{Standard error and confidence intervals}

Standard deviation describes variation among individual observations. \term{Standard error} describes uncertainty in an estimate. For a sample mean,
\[
  \operatorname{SE}(\bar{x})=\frac{s}{\sqrt{n}}.
\]
Standard error falls as $1/\sqrt{n}$: halving it requires about four times the data.

For independent, identically distributed observations with finite variance, sample means approach normality as sample size grows. Small samples, heavy tails, or dependence can invalidate the approximation. Under suitable conditions:
\[
  \bar{x}\pm z_{1-\alpha/2}\frac{s}{\sqrt{n}}.
\]

A valid 95\% confidence procedure covers the true parameter in about 95\% of repeated samples.

\subsection{The bootstrap}

The \term{bootstrap} estimates sampling variability by repeatedly resampling observations with replacement and recalculating the statistic.

\begin{mlsteps}
  \item Draw $n$ rows with replacement from the $n$ observed rows.
  \item Recalculate the statistic, such as a median or a difference in model loss.
  \item Repeat many times and use suitable quantiles as an interval.
\end{mlsteps}

Resampling must preserve dependence: use blocks for time series, groups for repeated measurements, and matched observations for paired comparisons.

\section{Probability basics}

A probability lies between 0 and 1. For events $A$ and $B$:

\begin{itemize}
  \item $P(A)$ is the probability of $A$.
  \item $P(A\cap B)$ is the probability that both occur.
  \item $P(A\cup B)$ is the probability that at least one occurs.
  \item $P(A\given B)$ is the probability of $A$ when $B$ is known.
\end{itemize}

The multiplication rule is
\[
  P(A\cap B)=P(A)P(B\given A).
\]
Events are independent when knowing one does not change the probability of the other:
\[
  P(A\cap B)=P(A)P(B).
\]

\section{Bayes' theorem}

Bayes' theorem reverses a conditional probability:
\[
  P(A\given B)=\frac{P(B\given A)P(A)}{P(B)}.
\]

\begin{description}[leftmargin=1.15in,style=nextline]
  \item[Prior] $P(A)$: belief before seeing new evidence.
  \item[Likelihood] $P(B\given A)$: how likely the evidence is if $A$ is true.
  \item[Posterior] $P(A\given B)$: updated belief after seeing the evidence.
\end{description}

\begin{workedexample}
A condition affects 1\% of people. A test detects 90\% of true cases and gives a false positive for 5\% of people without the condition. Then
\[
P(+)=0.90(0.01)+0.05(0.99)=0.0585.
\]
Therefore,
\[
P(\text{condition}\given +)=\frac{0.90(0.01)}{0.0585}\approx0.154.
\]
A positive result corresponds to about a 15.4\% chance of the condition, not 90\%, because the condition is rare.
\end{workedexample}

\begin{keyidea}
Always consider the base rate. Strong-looking evidence can still lead to a modest posterior probability when the event is rare.
\end{keyidea}

\section{Random variables and common distributions}

A random variable assigns a number to an uncertain outcome. Its expectation is a probability-weighted average. For a discrete random variable,
\[
  \mathbb{E}[X]=\sum_x xP(X=x).
\]

Common distributions include:

\begin{description}[leftmargin=1.25in,style=nextline]
  \item[Bernoulli] One binary outcome with probability $p$.
  \item[Binomial] Number of successes in $n$ independent Bernoulli trials.
  \item[Categorical] One result from several classes.
  \item[Poisson] Count of events in a fixed interval under simplifying assumptions.
  \item[Uniform] All values in a range have equal density, or all finite outcomes have equal probability.
  \item[Normal] Continuous values concentrated symmetrically around a mean.
\end{description}

Probability distributions connect directly to loss functions. Gaussian error leads naturally to squared-error regression, while a Bernoulli model leads naturally to binary cross-entropy.

\section{Covariance and correlation}
\label{sec:correlation}

Covariance describes whether two variables tend to move together:
\[
  \operatorname{Cov}(X,Y)=\mathbb{E}[(X-\mu_X)(Y-\mu_Y)].
\]
Pearson correlation standardizes covariance:
\[
  \rho_{X,Y}=\frac{\operatorname{Cov}(X,Y)}{\sigma_X\sigma_Y}.
\]
It lies between $-1$ and $1$ and measures linear association.

\begin{pitfall}
Correlation does not establish causation: confounding, reverse direction, or selection may explain it.
\end{pitfall}

A correlation near zero also does not rule out a curved or other nonlinear relationship.

\begin{optionaltopic}
\textbf{Likelihood and maximum likelihood estimation.} Given data $x_1,\ldots,x_n$, the likelihood treats the probability of the observed data as a function of a parameter $\theta$:
\[
L(\theta)=\prod_{i=1}^{n}p(x_i\given\theta).
\]
Maximum likelihood selects the parameter making the observations most probable. Taking logs converts products into sums:
\[
\hat{\theta}_{\mathrm{MLE}}=\argmax_{\theta}\sum_{i=1}^{n}\log p(x_i\given\theta).
\]
Many familiar machine-learning losses are negative log-likelihoods.
\end{optionaltopic}

\begin{optionaltopic}
\textbf{Hypothesis tests.} A p-value is the probability, under the null model, of a result at least as extreme as observed---not the probability the null is true. Report effect size and uncertainty too.
\end{optionaltopic}

\begin{optionaltopic}
\textbf{Entropy} measures uncertainty; \textbf{cross-entropy} evaluates predicted probabilities; \textbf{KL divergence} compares distributions (\chapref{ch:classification}, \chapref{ch:evaluation}).
\end{optionaltopic}

\section{Scalars, vectors, and matrices}

\begin{description}[leftmargin=1.1in,style=nextline]
  \item[Scalar] One number, such as temperature $23$.
  \item[Vector] An ordered list of numbers, such as one house $\bm{x}=[120,3,8]^{\mathsf T}$.
  \item[Matrix] A rectangular table of numbers, such as a dataset with one row per house and one column per feature.
\end{description}

For $n$ observations and $p$ features, the design matrix has shape
\[
  X\in\mathbb{R}^{n\times p}.
\]
The first number is the number of rows; the second is the number of columns.

\begin{workedexample}
Five hundred students with four features give $X\in\mathbb{R}^{500\times4}$. Lab~1's Student Performance matrix has 649 rows and 32 feature columns.
\end{workedexample}

\section{Matrix shapes: the rule that governs everything}

For matrix multiplication, inner dimensions must match; outer dimensions determine the result.

\begin{mlfigure}
\begin{tikzpicture}[font=\small]
\fill[MLPaleBlue,draw=MLBlue,rounded corners=1mm] (0,0) rectangle (2.2,1.9);
\node[font=\bfseries\color{MLNavy}] at (1.1,0.95) {$X$};
\node[font=\scriptsize\color{MLMuted}] at (1.1,-0.35) {$(n \times p)$};
\node[font=\Large\color{MLInk}] at (2.75,0.95) {$\times$};
\fill[MLPaleTeal,draw=MLTeal,rounded corners=1mm] (3.3,0.35) rectangle (4.3,1.55);
\node[font=\bfseries\color{MLNavy}] at (3.8,0.95) {$\bm{w}$};
\node[font=\scriptsize\color{MLMuted}] at (3.8,-0.35) {$(p \times 1)$};
\node[font=\Large\color{MLInk}] at (4.85,0.95) {$=$};
\fill[MLPaleOrange,draw=MLOrange,rounded corners=1mm] (5.4,0) rectangle (6.1,1.9);
\node[font=\bfseries\color{MLNavy}] at (5.75,0.95) {$\hat{\bm{y}}$};
\node[font=\scriptsize\color{MLMuted}] at (5.75,-0.35) {$(n \times 1)$};
\draw[{Stealth[length=2mm]}-{Stealth[length=2mm]},thick,MLRed] (1.75,-0.75) -- (3.35,-0.75);
\node[font=\scriptsize\color{MLRed}] at (2.55,-1.05) {these must match};
\draw[{Stealth[length=2mm]}-,thick,MLGreen] (0.4,2.25) -- (0.4,2.0);
\draw[-{Stealth[length=2mm]},thick,MLGreen] (0.4,2.25) -- (5.6,2.25);
\draw[-{Stealth[length=2mm]},thick,MLGreen] (5.6,2.25) -- (5.6,2.0);
\node[font=\scriptsize\color{MLGreen}] at (3.0,2.5) {these survive into the answer};
\end{tikzpicture}
\end{mlfigure}

\section{Vector arithmetic}

Vectors of the same length can be added component by component:
\[
\begin{bmatrix}1\\2\end{bmatrix}+
\begin{bmatrix}3\\4\end{bmatrix}=
\begin{bmatrix}4\\6\end{bmatrix}.
\]
Multiplying by a scalar stretches or reverses a vector:
\[
3\begin{bmatrix}1\\2\end{bmatrix}=
\begin{bmatrix}3\\6\end{bmatrix}.
\]

\section{Dot products: weighted sums and alignment}

The dot product of two vectors is
\[
  \bm{a}^{\mathsf T}\bm{b}=\sum_{j=1}^{p}a_jb_j.
\]
It has two useful interpretations.

\begin{enumerate}
  \item It creates a weighted sum. Linear models use $\bm{w}^{\mathsf T}\bm{x}+b$.
  \item It measures alignment. A large positive dot product means the vectors point in similar directions; a negative one means they point in opposing directions.
\end{enumerate}

\begin{workedexample}
For $\bm{w}=[2,-1]^{\mathsf T}$ and $\bm{x}=[3,4]^{\mathsf T}$,
\[
\bm{w}^{\mathsf T}\bm{x}=2(3)+(-1)(4)=2.
\]
A linear model with bias $b=0.5$ would predict $2.5$.
\end{workedexample}

\section{Norms: measuring vector size}

The $L_1$ norm is the sum of absolute values:
\[
  \|\bm{x}\|_1=\sum_j|x_j|.
\]
The $L_2$ norm is the ordinary straight-line length:
\[
  \|\bm{x}\|_2=\sqrt{\sum_jx_j^2}.
\]
These norms appear in distance calculations and regularization.

\section{Distance and similarity}

A model often needs to decide whether two observations are close or similar.

\begin{mltable}
\begin{tabularx}{\linewidth}{@{}>{\bfseries\leavevmode\color{MLNavy}}p{0.20\linewidth}>{\raggedright\arraybackslash}p{0.28\linewidth}>{\raggedright\arraybackslash}X@{}}
\toprule
\rowcolor{MLPaleBlue!92}
Measure & Formula or idea & Useful when \cr
\midrule
Euclidean distance & $\|\bm{x}-\bm{z}\|_2$ & ordinary geometric distance; features are scaled \cr
Manhattan distance & $\|\bm{x}-\bm{z}\|_1$ & movement along grid-like axes; more robustness to a single large difference \cr
Cosine similarity & $\frac{\bm{x}^{\mathsf T}\bm{z}}{\|\bm{x}\|_2\|\bm{z}\|_2}$ & direction matters more than magnitude, such as word-count vectors \cr
Jaccard similarity & overlap divided by union & sets or binary presence/absence features \cr
Hamming distance & number of unequal positions & binary strings or categorical codes \cr
\bottomrule
\end{tabularx}
\end{mltable}

\begin{pitfall}
Distance-based models can be meaningless when features use very different scales. Standardize or otherwise justify the scale before calculating distances.
\end{pitfall}

\subsection*{The same two points can look different under different measures}

Euclidean, Manhattan, Jaccard, Minkowski, and cosine measures encode different meanings of ``close''. Choose one for the representation and task.

\begin{workedexample}
For points $P=(1,2)$ and $Q=(4,6)$,
\[
 d_2(P,Q)=\sqrt{(4-1)^2+(6-2)^2}=5,
\]
while
\[
 d_1(P,Q)=|4-1|+|6-2|=7.
\]
Euclidean distance measures the straight-line path; Manhattan distance measures the axis-by-axis path.
\end{workedexample}

The \term{Minkowski distance} unifies both:
\[
 d_p(\bm{x},\bm{z})=
 \left(\sum_j |x_j-z_j|^p\right)^{1/p}.
\]
Minkowski uses Manhattan distance at $p=1$ and Euclidean distance at $p=2$. Choose $p$ to match the task.

\begin{workedexample}
For sets $A=\{0,1,2,5,6\}$ and $B=\{0,2,3,4,5,7,9\}$, the intersection has three elements and the union has nine. Therefore
\[
 J(A,B)=\frac{|A\cap B|}{|A\cup B|}=\frac{3}{9}=0.33.
\]
Jaccard measures overlap among present items, ignoring shared absences.
\end{workedexample}

\begin{keyidea}
Choose representation and distance together; a computable distance need not be meaningful.
\end{keyidea}

\section{Cosine similarity visually}

Cosine similarity depends on the angle $\theta$ between two vectors:
\[
  \cos(\theta)=\frac{\bm{x}^{\mathsf T}\bm{z}}{\|\bm{x}\|_2\|\bm{z}\|_2}.
\]

\begin{mlfigure}
\begin{tikzpicture}[scale=1.05]
\draw[-{Stealth[length=2.5mm]},thick,MLBlue] (0,0)--(4,1.2) node[right] {$\bm{x}$};
\draw[-{Stealth[length=2.5mm]},thick,MLTeal] (0,0)--(3,2.7) node[above] {$\bm{z}$};
\draw[MLOrange,thick] (1.1,0.33) arc[start angle=16.7,end angle=42, radius=1.15];
\node[MLOrange] at (1.45,0.82) {$\theta$};
\draw[->,MLMuted] (-0.3,0)--(4.5,0);
\draw[->,MLMuted] (0,-0.3)--(0,3.2);
\end{tikzpicture}
\end{mlfigure}

A cosine similarity near 1 means similar direction, near 0 means roughly perpendicular, and near $-1$ means opposite direction.

\section{Matrices and linear transformations}

A matrix can transform one vector into another:
\[
  \bm{y}=A\bm{x}.
\]
Depending on its values, $A$ can rotate, stretch, shrink, reflect, or combine coordinates. A linear regression model applies this same idea to many rows at once:
\[
  \hat{\bm{y}}=X\bm{\beta}.
\]

\begin{intuition}
Matrix multiplication forms weighted combinations; compatible shapes determine which products exist.
\end{intuition}

\section{Transpose, identity, and inverse}

The transpose switches rows and columns. If $A$ is $m\times p$, then $A^{\mathsf T}$ is $p\times m$.

The identity matrix $I$ acts like the number 1:
\[
  I\bm{x}=\bm{x}.
\]
When an inverse exists,
\[
  A^{-1}A=I.
\]
Solve linear systems numerically rather than explicitly computing inverses.

\section{Gradients: which way changes the loss?}

For a function with several parameters, the gradient collects all partial derivatives:
\[
  \nabla_{\bm{\theta}}J=
  \begin{bmatrix}
  \frac{\partial J}{\partial\theta_1}\\
  \vdots\\
  \frac{\partial J}{\partial\theta_p}
  \end{bmatrix}.
\]
The gradient points toward the fastest increase in $J$. To reduce a loss, gradient descent moves in the opposite direction:
\[
  \bm{\theta}\leftarrow\bm{\theta}-\alpha\nabla_{\bm{\theta}}J.
\]

For step sizes and a worked gradient-descent update, see \secref{sec:gradient-descent}.

\begin{optionaltopic}
\textbf{Eigenvectors and SVD.} An eigenvector is scaled by a square matrix. SVD factors a matrix into rotations/reflections and scaling. Both reveal directions used in PCA (\chapref{ch:dimred}).
\end{optionaltopic}

\begin{optionaltopic}
\textbf{Hessian.} This matrix of second derivatives describes curvature: positive in bowl-like directions, negative in hill-like directions. It supports advanced optimization.
\end{optionaltopic}

\begin{keyidea}
Vectors represent observations and parameters; matrices represent datasets and transformations; dot products create weighted sums; distances compare observations; and gradients tell optimization which direction to move.
\end{keyidea}

\begin{quickcheck}
\begin{enumerate}
  \item When would the median and IQR describe a distribution more faithfully than the mean and standard deviation?
  \item If \(X\) has shape \(n\times p\), what shape must a coefficient vector have for \(X\bm{w}\) to produce one prediction per row?
  \item Why can standardizing features change a nearest-neighbor or clustering result even though the observations themselves did not change?
\end{enumerate}
\end{quickcheck}
